%% The measured results of the rescaling (use A) and of the blended projection extension (use B),
%% 2026-10-31. Numbers from the curated tables
%% docs/gradient-controlled-level-set/gcls-level-set-article/data/tables/rescale_*.csv and
%% projectionBlend_*.csv (static_source_metrics.py and projection_blend_metrics.py at 9d483075,
%% Lichtenberg job 55814177; runs at 049cfa66). STATUS.md 11.28, items 7 to 9.
\section{The rescaling and the blended extension: measured}
\label{sec:rescaling-results}

\paragraph{Set-up.} The static circle of this report, uniform hexahedra, $N=64$ to 512,
\texttt{REDISTANCER rescaling}, no source, Euler in pseudo-time (Section~\ref{sec:rescaling-fvm},
item 3), the band of the sign-change cells and four layers. ``Once'': one event at the first
step, read at $\tau=1$. ``Every step'': one event per step to $\tau=10$. The error is
$\mathrm{dist3}_2$, the $L_2$ norm of $(\psi-d)/h$ over $|d|\le3h$, so order $p$ is a pointwise
$O(h^{p+1})$. The pre-registered lines R1 to R6 are in the headers of
\path{config/staticSource/staticSourceRescale*2D.yaml}.

\begin{table}[h]\centering\small
\caption{One event: $\mathrm{dist3}_2$ after the event, and its order over $N=64$ to 512.
\Meas}
\label{tab:rescale-once}
\begin{tabular}{llllll}
\toprule
$|\nabla\psi|$ & map & start & $N=64$ & $N=512$ & order \\
\midrule
centred & second (0b) & $r^2-R^2$ & $4.9\times10^{-15}$ & $4.5\times10^{-14}$ & round-off \\
centred & second (0b) & blend 0.1 & $8.5\times10^{-4}$ & $1.5\times10^{-5}$ & 1.96 \\
centred & second (0b) & $d$ & $9.4\times10^{-4}$ & $1.6\times10^{-5}$ & 1.96 \\
centred & first (0a) & $r^2-R^2$ & $1.2\times10^{-1}$ & $1.6\times10^{-2}$ & 1.00 \\
centred & first (0a) & blend 0.1 & $1.3\times10^{-2}$ & $1.6\times10^{-3}$ & 1.00 \\
one-sided & first (0a) & $d$ & $2.6\times10^{-2}$ & $3.3\times10^{-3}$ & 0.99 \\
one-sided & second (0b) & $r^2-R^2$ & $6.9\times10^{-2}$ & $8.7\times10^{-3}$ & 1.00 \\
\bottomrule
\end{tabular}
\end{table}

\paragraph{One event (\Meas).}
\begin{enumerate}
\item The centred second-order map gives the orders that Section~\ref{sec:rescaling-fvm} derives
for uniform hexahedra: exact for $r^2-R^2$, order 2 otherwise. The zero crossings and the DA
volume after the event equal those of $d$ (R2, R3: relative difference 0.000 at $N\ge128$), and
$d$ is a fixed point to $1.6\times10^{-5}h$ at $N=512$ (R4). The line R1 for $r^2-R^2$ reads
FAIL as written: its least-squares order of a round-off error is $-1.06$, and the read-out had no
clause for that case. The verdict stays as written; the value is round-off.
\item The centred first-order map is order 1 where the model term $-(B/2A)d^2$ is present (R1
PASS). On $(r^2-R^2)/(2R)$ it does not reduce the error: the model term has the size of the start
error and the opposite sign.
\item The one-sided (Rouy--Tourin) $|\nabla\psi|$ is a first-order estimate, so every map with
it is order 1: it moves $d$ itself (R4 FAIL), and with the second-order map it moves the zero
crossings by 0.9 to 1.8 times the error of $d$ and changes the DA volume (R2, R3 FAIL).
\end{enumerate}

\paragraph{Every step (\Meas).} Every arm FAILS R5.
\begin{enumerate}
\item Centred, first and second order, least-squares and Gauss Hessian: unstable from the first
events. From $d$ at $N=64$ (second order) the band mean of $|\nabla\psi|$ is 37.6 after 13
events and $2.7\times10^4$ after 26. At $\tau=10$, $\mathrm{dist3}_2$ is $3\times10^4$ to
$6\times10^6$ for every start and every $N$; the zero crossings moved $0.3h$ to $0.7h$, the DA
volume grew 10 to 15\,\%.
\item One-sided, first order: bounded. It settles at $\mathrm{dist3}_2=1.13$ on every rung (order
0); the zero crossings moved $0.24h$ to $0.41h$; the volume changed by $-3.3\,\%$ ($N=64$) to
$-0.2\,\%$ ($N=512$). One-sided, second order: $\mathrm{dist3}_2$ between 3 and 41, not
monotone in $N$ or $\tau$.
\item R6 reads PASS as written (Gauss over least-squares Hessian at $\tau=10$: 0.89 to 1.02) only
because both arms diverged to the same level.
\end{enumerate}

\paragraph{Does the noise of the FVM Hessian destroy the solution?} On uniform hexahedra the
question does not arise: the first-order map, which uses no Hessian, diverges in the same way, and
the Gauss Hessian gives the same result as the least-squares one. The centred gradient destroys
the solution first, by the mechanism of Section~\ref{sec:rescaling-fvm}, item 2 (one application
doubles the mode $(-1)^{i+j}d$). The polyhedral prediction (item 4 there) is not measured.

\paragraph{Assessment (use A on the UFVM line).} The rescaling is a correct ONE-TIME map, not an
alternative to a source term with these gradients. Applied once (centred, second order) it gives a
second-order distance and keeps the zero set and the DA volume. Applied every step it diverges
(centred) or keeps an $O(h)$ error and moves the interface (one-sided). Open: an event every $k$
steps, $k>1$, which bounds the growth by $2^{n/k}$ only if the transport between events damps the
mode; and the semi-Lagrangian version with the quadratic fit (task T6 of the SL session), whose
fit is not blind to the odd-even mode (Section~\ref{sec:rescaling-sl}).

\paragraph{Use B: the blended projection extension (\Meas).} \texttt{leiaLevelSetFoam}, Eulerian
transport, reversed flows: 2D vortex $N=64,128,256$ on 4 ranks, 3D shear $N=32,64,128$ on 16
ranks. The 3D deformation study had a setup error (the case template had no \texttt{law} block,
every \texttt{gradientControl} arm stopped at step 0); it is void and runs again.

\begin{table}[h]\centering\small
\caption{Extension (second distance) against no extension, finest rung; ratio extension over
baseline, and the observed orders. Gradient error in the band at $T/2$, shape at $T$, volume at
$T$. \Meas}
\label{tab:pb}
\begin{tabular}{lllll}
\toprule
case & gradient ratio & shape ratio & shape order & volume at $T$ \\
\midrule
2D vortex, $N=256$ & 0.46 & 5.1 & 2.53 (base 3.23) & $1.5\times10^{-3}$ (base $1.9\times10^{-4}$) \\
3D shear, $N=128$ & 0.43 & 1.56 & 1.28 (base 1.87) & $5.1\times10^{-2}$ (base $4.7\times10^{-3}$) \\
\bottomrule
\end{tabular}
\end{table}

\begin{enumerate}
\item B1 FAILS on both cases: the band gradient error halves, but the shape and the volume at $T$
are worse and converge more slowly (Table~\ref{tab:pb}).
\item The band gradient error does not converge in any arm, the baseline included (orders
$-0.3$ to 0). The band mean of $|\nabla\psi|$ at $T/2$ is 2.4 (2D) and 3.2 (3D) without the
extension and 1.45 to 1.9 with it: the extension reduces the steepening, but does not hold
$|\nabla\psi|=1$.
\item With the small source (\texttt{linearQ}, $\mu=0.1$) the gradient error falls further, to
0.77 to 0.86 of the extension alone (B3 PASS as written), and the shape and the volume become
worse again.
\item The null pairs are bit-identical and every case completed (B4 PASS).
\item \Hyp{} The baseline's high shape order at $T$ comes from the time reversal of a linear
transport with a prescribed velocity; an extension that depends on $\psi$ breaks that
cancellation. In 2D the extension reduces the volume error at $T/2$ 16 times, which points the
same way. A flow that does not reverse separates the two.
\end{enumerate}
